\subsection{Propuesta de Solución}

El problema central identificado en el centro odontológico
Bellidental radica en la fragmentación del flujo de información
clínico-administrativa, sostenido por registros manuales dispersos
en herramientas desconectadas entre sí. Ante esta situación, se
propone el desarrollo de un sistema de información centralizado
que unifique en un único entorno los procesos de registro de la
historia clínica, generación de recetas, agendamiento de citas y
registro de pagos. La solución no consiste en reemplazar una
herramienta por otra de forma aislada, sino en estructurar un
flujo de información continuo donde cada acción registrada quede
automáticamente vinculada al resto de los datos del paciente,
eliminando la necesidad de trasvasar información de forma manual
entre plataformas separadas. La variable independiente de la
solución es el nivel de estandarización e integración del flujo de
información clínico-administrativa; su efecto será medido a través
de la variable dependiente: el tiempo promedio que el personal
dedica a completar el ciclo administrativo de cada paciente,
expresado en minutos por ciclo completo de atención.

El primer paso de la propuesta consiste en documentar con rigor el
estado actual del centro. Para ello, se realizarán entrevistas al
personal directamente involucrado en la atención y se observará de
forma directa el desarrollo de las jornadas de trabajo. Esta fase
permitirá identificar qué procesos existen actualmente, en qué
puntos se producen interrupciones o demoras, qué información se
pierde entre una herramienta y otra, y qué exigencias normativas
del Ministerio de Salud Pública debe cumplir el sistema, incluyendo
los formatos y campos del Formulario 033. Los hallazgos obtenidos
constituirán la base objetiva sobre la cual se tomará cada
decisión de diseño posterior, evitando que el sistema responda a
supuestos del desarrollador en lugar de a las necesidades reales
de la clínica.

El segundo paso consiste en diseñar el modelo de información que
estructurará el sistema. A partir del diagnóstico previo se
definirán las reglas de funcionamiento de cada proceso, la forma
en que los datos del paciente se organizarán y circularán dentro
del sistema, y los mecanismos que garanticen el cumplimiento de la
normativa vigente en materia de historias clínicas y
confidencialidad de datos. Esta etapa producirá un modelo
estructurado, revisado y validado antes de iniciar cualquier
desarrollo técnico. El resultado concreto de esta fase será la
especificación de los módulos funcionales del sistema y la
descripción del flujo unificado que conectará el acto clínico con
el registro administrativo y el control financiero.

El tercer paso comprende la construcción e instalación del sistema
en el entorno real de trabajo de Bellidental. Una vez en
funcionamiento, se aplicará el diseño de evaluación preexperimental
de medición antes y después: primero se habrá registrado el tiempo
promedio del ciclo administrativo por paciente bajo los métodos
tradicionales; luego, con el sistema en operación, se realizará
la misma medición bajo las mismas condiciones y criterios. La
comparación de ambas mediciones, apoyada en el cálculo del
promedio, la desviación estándar y la prueba $t$ de Student para
muestras relacionadas, permitirá determinar con datos numéricos
concretos si el sistema logra reducir la carga administrativa y en
qué magnitud. Los ciclos de atención a medir en cada etapa serán
seleccionados de forma aleatoria para garantizar representatividad
y evitar sesgos en la comparación.

La viabilidad técnica de la propuesta se sustenta en que la
clínica ya opera con dispositivos digitales de uso cotidiano y
dispone de conectividad básica en el consultorio, condiciones
coherentes con los requerimientos de infraestructura declarados en
la delimitación del estudio. El sistema será diseñado considerando
estas capacidades tecnológicas reales del entorno, evitando
dependencias de infraestructura compleja o especializada que
encarezcan el proyecto o dificulten su mantenimiento futuro.
Desde el punto de vista económico, la solución representa una
inversión concentrada en el período de desarrollo —cuantificada en
la sección de recursos— tras la cual la clínica podrá reducir
costos recurrentes asociados a la impresión de formularios, la
pérdida de citas por ausencia de confirmación oportuna y el tiempo
del personal destinado a tareas manuales repetitivas que hoy
consumen minutos productivos de cada jornada.

La viabilidad operativa se fundamenta en que el sistema será
construido a partir del flujo de trabajo ya documentado del propio
personal, lo que reduce la resistencia al cambio y facilita la
adopción en el entorno cotidiano. Al estar diseñado sobre los
procesos reales de Bellidental y no sobre un modelo genérico, el
personal reconocerá sus propios procedimientos reflejados en el
sistema, lo que acorta la curva de aprendizaje. Desde una
perspectiva social, la propuesta contribuye a mejorar la calidad
de atención percibida por el paciente, quien contará con mayor
claridad sobre su tratamiento, sus indicaciones terapéuticas y el
estado de su cuenta, información que actualmente depende de que
el personal recuerde comunicarla de forma manual en cada consulta.